\chapter{Étude de l'existant et positionnement}
\label{ch:etat-art}

\section{Périmètre et méthode}
\label{sec:etat-methode}

Cette étude positionne LexHub face aux offres accessibles à un cabinet marocain. La grille est stable : atelier voix~$\rightarrow$~plaidoirie, gestion de cabinet, SaaS multi-tenant, travail en équipe, et \textbf{maîtrise des faits}.

\paragraph{Méthode AS-IS / TO-BE.}
\begin{itemize}[leftmargin=1.4em]
  \item \textbf{AS-IS} --- ce que le cabinet peut acheter ou assembler aujourd'hui ;
  \item \textbf{Écart} --- ce qui manque pour soutenir \textbf{dossier $\rightarrow$ pièces $\rightarrow$ rédaction $\rightarrow$ livrable} ;
  \item \textbf{TO-BE} --- le SI de cabinet LexHub, où l'IA est un \textbf{service}, pas le produit.
\end{itemize}

La même grille structure l'urbanisation interne du SI (section~\ref{sec:madi}). Sources : \texttt{smartavocat.ma}, \texttt{avocato.ma}, \texttt{avocatpro.ma}, \texttt{legaliq.ma}, \texttt{kanoun.com}, \texttt{adala.ai}, \texttt{diskod.ma}.

\section{AS-IS --- paysage concurrentiel}
\label{sec:as-is-marche}

Quatre familles coexistent, souvent \textbf{sans pont} : gestion de cabinet ; dictée bureautique ; assistant IA juridique ; assemblage manuel (Word + Drive + messagerie).

\begin{table}[H]
\centering
\caption{Benchmarking marché --- AS-IS vs besoin LexHub}
\label{tab:benchmark}
\footnotesize
\renewcommand{\arraystretch}{1.35}
\begin{tabular}{@{}p{2.4cm}p{4.6cm}p{5.6cm}@{}}
\toprule
\textbf{Offre} & \textbf{Services offerts (AS-IS)} & \textbf{Écart face à LexHub} \\
\midrule
SmartAvocat \newline {\scriptsize smartavocat.ma} &
Gestion cabinet, dossiers bilingues &
Pas de flux plaidoirie vocale \\
\midrule
Avocato \newline {\scriptsize avocato.ma} &
Gestion + audiences / portail justice &
Pas de parcours dictée $\rightarrow$ document métier \\
\midrule
AvocatPro \newline {\scriptsize avocatpro.ma} &
Gestion + calendrier judiciaire &
Pas de cœur plaidoirie \\
\midrule
Legaliq \newline {\scriptsize legaliq.ma} &
IA juridique / fiscale, documents &
Pas un SaaS cabinet complet \\
\midrule
Kanoun \newline {\scriptsize kanoun.com} &
IA droit marocain + OHADA &
Pas multi-tenant cabinet \\
\midrule
Adala \newline {\scriptsize adala.ai} &
Assistant IA, Q\&A, documents &
Pas d'équipe cabinet \\
\midrule
Dragon Legal / Diskod \newline {\scriptsize diskod.ma} &
Dictée desktop (Nuance) &
Pas SaaS, pas dossiers LexHub \\
\midrule
\textbf{LexHub} \newline {\scriptsize lexhub.space} &
Atelier 5 étapes ; dossiers ; équipe ; SaaS &
--- \\
\bottomrule
\end{tabular}
\end{table}

\begin{table}[H]
\centering
\caption{Lecture urbanisée du marché AS-IS}
\label{tab:as-is-lecture}
\small
\begin{tabularx}{\textwidth}{@{}>{\bfseries}p{2.8cm}XX@{}}
\toprule
Couche & Constat marché & Conséquence pour le cabinet \\
\midrule
Métier & Processus découpé (gestion \emph{ou} dictée \emph{ou} chat) & Ressaisie, perte de contexte \\
Applications & Outils en silos & Chaîne de valeur brisée \\
Données & Word / Drive souvent maîtres de fait & Pas de cadastre unique du dossier \\
Infrastructure & Poste local + cloud grand public & Secret professionnel difficile à garantir \\
\bottomrule
\end{tabularx}
\end{table}

\section{Écart --- ce que le marché ne couvre pas}
\label{sec:gap-benchmark}

\begin{itemize}[leftmargin=1.4em]
  \item pas de chaîne unique \textbf{dossier $\rightarrow$ pièces $\rightarrow$ dictée $\rightarrow$ acte} ;
  \item pas de rôle assistant borné aux plaidoiries \textbf{validées} ;
  \item pas de règle « LLM $\neq$ maître des faits » ;
  \item pas de SaaS multi-cabinet couvrant gestion \emph{et} production juridique.
\end{itemize}

\section{TO-BE --- positionnement LexHub}
\label{sec:to-be-position}

\begin{table}[H]
\centering
\caption{Du marché AS-IS à la cible TO-BE LexHub}
\label{tab:as-to-be}
\small
\renewcommand{\arraystretch}{1.3}
\begin{tabularx}{\textwidth}{@{}>{\bfseries}p{2.6cm}XX@{}}
\toprule
Dimension & AS-IS marché & TO-BE LexHub \\
\midrule
Cœur métier & Gestion \emph{ou} dictée \emph{ou} chat & Dossier + plaidoirie (chaîne unique) \\
Rôle de l'IA & Produit affiché / chatbot & Service (STT, OCR, révision) \\
Données & Fichiers Word / Drive dispersés & Postgres (vérité) + MinIO (fichiers) \\
Équipe & Peu de validation de brouillons & Assistant = plaidoiries validées \\
Modèle & Licence poste / freemium IA & SaaS multi-tenant (plans, packs) \\
Urbanisation & Silos applicatifs & Blocs MADI (section~\ref{sec:madi}) \\
\bottomrule
\end{tabularx}
\end{table}

\begin{figure}[H]
\centering
\begin{tikzpicture}[
  box/.style={draw=royal!60,fill=royal!5,rounded corners=3pt,align=center,font=\small,inner sep=5pt,minimum height=1.35cm},
  mid/.style={draw=gold!80!black,fill=gold!12,rounded corners=3pt,align=center,font=\small\bfseries,inner sep=5pt},
  arr/.style={-{Stealth},thick,royal!75}
]
  \node[box,minimum width=3.6cm] (asis) at (0,0) {AS-IS\\Outils en silos};
  \node[mid,minimum width=3.2cm] (gap) at (5.0,0) {ÉCART\\Process non uni};
  \node[box,minimum width=3.8cm] (tobe) at (10.0,0) {TO-BE\\SI LexHub};
  \draw[arr] (asis) -- (gap);
  \draw[arr] (gap) -- (tobe);
\end{tikzpicture}
\caption{Chaîne méthodologique du benchmarking : AS-IS $\rightarrow$ écart $\rightarrow$ TO-BE}
\label{fig:as-to-be-chain}
\end{figure}

\section{Synthèse et choix retenus}

\begin{enumerate}[leftmargin=1.4em]
  \item Un \textbf{SI de cabinet}, pas un chatbot générique.
  \item Un atelier voix~$\rightarrow$~document en cinq étapes, distinct d'un logiciel de greffe.
  \item Fidélité à la dictée : aucun fait inventé.
  \item SaaS facturable (chapitre~\ref{ch:economie}) et urbanisation MADI (section~\ref{sec:madi}).
\end{enumerate}
